\scanpage{140}\mspage{- 135 -}\hfill IX

qui montre que $\Phi$ \uline{est canoniquement isomorphe au sous-groupe de torsion de} $H^1(I,U)$.

D'autre part, de la suite exacte (11.3.3) on déduit un isomorphisme, pour tout $n=\ell^r$ :
\[
 H^1(I,T_n)\simeq\mathrm{Hom}(I,T_n)\overset{\sim}{\longleftarrow}\mathrm{Hom}(T,T_n)\overset{\sim}{\longleftarrow}\mathbb Z/n\mathbb Z \tag*{(11.3.9)}
\]
(NB on fait opérer $I$ trivialement sur $T$, donc sur $T_n=T/nT$) ; d'où, pour tout $I$-$\mathbb Z_\ell$-Module fini $M$, si
\[
 M'=\mathrm{Hom}(M,T\otimes D_\ell) \tag*{(11.3.10)}
\]
désigne son ``dual de Cartier'', un accouplement canonique
\[
 H^i(I,M)\otimes H^{1-i}(I,M')\longrightarrow H^1(I,T\otimes D_\ell)\simeq D_\ell\ ,\tag*{(11.3.11)}
\]
où le dernier isomorphisme est déduit des isomorphismes (11.3.9) par passage à la limite inductive sur les puissances $n$ de $\ell$. Il est bien connu que ces accouplements sont des \uline{dualités parfaites} (cf. par exemple SGA 5 I 5). Prenant dans ces accouplements $M$ de la forme $U_n$, donc $M'=U'_n$, on trouve par passage à la limite un accouplement parfait
\[
 H^{\unsure{o|0}}(I,U\otimes D_\ell)\times H^1(I,U')\longrightarrow D_\ell\ ,\tag*{(11.3.12)}
\]
où le deuxième facteur du premier membre est un $\mathbb Z_\ell$-module de type fini. Un nouveau passage à la limite dans les accouplements (11.3.12) (pour des morphismes de transition de la forme $\ell^r.\mathrm{id}$) nous donne un deuxième accouplement parfait
\[
 H^o(I,U\otimes\mathbb Q_\ell)\times H^1(I,U'\otimes\mathbb Q_\ell)\longrightarrow D_\ell \tag*{(11.3.13)}
\]
de vectoriels de dimension finie sur $\mathbb Q_\ell$. Désignant par $i'$, $j'$ les
% archive sheet number: 139
% stamp: numbered label 453. Three broken vertical impressions recur at the left, without an identifiable text line.
